% GENERATED FILE. Edit canonical lessons, then run npm run generate:books.
% canonical-source-hash: fnv1a64:269ae99c8eb79ed4
% canonical-lessons: GU-C406-vijha, GU-C406-sima, GU-C406-jakat, GU-C406-phi, GU-C406-thalavvu

\chapter{At the Border}
\label{ch:gu-at-the-border}
\hlchaptermodality{406}

\begin{chapteropening}
\textbf{By the end of this chapter:} \emph{I can get across a border with a visa and my luggage.}

Complete the last lesson of chapter 406: I can get across a border with a visa and my luggage.
\end{chapteropening}

\section[\texorpdfstring{\gu{વિઝા} (vijhā)}{વિઝા (vijhā)}]{\gu{વિઝા} (vijhā) — a visa}

\label{lesson:GU-C406-vijha}



\begin{quote}
Before the new one: say the Gujarati for a suitcase, then the Gujarati for a circle in the air.
\end{quote}

\subsection*{You'll want to know: \gu{વિઝા}}
\textbf{\gu{વિઝા}} — \emph{vijhā} — \textquotedblleft{}a visa\textquotedblright{}.

\begin{quote}
\gu{પાસપોર્ટ} \gu{સાથે} \gu{વિઝા} \gu{પણ} \gu{જોઈએ}.
\end{quote}

\begin{quote}
\emph{pāsporṭ sāthe vijhā paṇ joīe.}
\end{quote}

\begin{quote}
As well as a passport you need a visa.
\end{quote}

A masculine noun.

\begin{cousinweb}
\textbf{\gu{વિઝા}} (\emph{vijhā}) is borrowed from English \emph{visa}, the stamp or paper that lets you into another country.
\end{cousinweb}

\begin{grammarlens}[title={What you've built}]
One more word for getting around.
\end{grammarlens}

\subsection*{Your turn}
\begin{itemize}
\raggedright
  \item \emph{Say it:} \emph{vijhā}
  \item \emph{Say it:} \emph{vijhā}, once more
  \item \emph{Say it:} \emph{pāsporṭ sāthe vijhā paṇ joīe.} — As well as a passport you need a visa.
  \item \emph{Recall:} say the Gujarati for a suitcase, then the Gujarati for a circle in the air, then say \emph{vijhā} again
\end{itemize}

\begin{tcolorbox}[breakable,colback=teal!4,colframe=teal!35!black,title={Before you move on}]
What does \gu{વિઝા} mean? (\textquotedblleft{}A visa\textquotedblright{}.) Say it once more.
\end{tcolorbox}
\section[\texorpdfstring{\gu{સીમા} (sīmā)}{સીમા (sīmā)}]{\gu{સીમા} (sīmā) — a border; a boundary}

\label{lesson:GU-C406-sima}



\begin{quote}
Before the new one: say the Gujarati for a circle in the air, then the Gujarati for a visa.
\end{quote}

\subsection*{You'll want to know: \gu{સીમા}}
\textbf{\gu{સીમા}} — \emph{sīmā} — \textquotedblleft{}a border; a boundary\textquotedblright{}.

\begin{quote}
\gu{આ} \gu{ગામ} \gu{સીમાની} \gu{નજીક} \gu{છે}.
\end{quote}

\begin{quote}
\emph{ā gām sīmānī najīk che.}
\end{quote}

\begin{quote}
This village is near the border.
\end{quote}

A feminine noun.

\begin{cousinweb}
\textbf{\gu{સીમા}} (\emph{sīmā}) is a tatsama, Sanskrit \emph{sīmā}, \textquotedblleft{}a boundary, a border\textquotedblright{} (Monier-Williams): the line where one land ends and the next begins.
\end{cousinweb}

\begin{grammarlens}[title={What you've built}]
One more word for getting around.
\end{grammarlens}

\subsection*{Your turn}
\begin{itemize}
\raggedright
  \item \emph{Say it:} \emph{sīmā}
  \item \emph{Say it:} \emph{sīmā}, once more
  \item \emph{Say it:} \emph{ā gām sīmānī najīk che.} — This village is near the border.
  \item \emph{Recall:} say the Gujarati for a circle in the air, then the Gujarati for a visa, then say \emph{sīmā} again
\end{itemize}

\begin{tcolorbox}[breakable,colback=teal!4,colframe=teal!35!black,title={Before you move on}]
What does \gu{સીમા} mean? (\textquotedblleft{}A border; a boundary\textquotedblright{}.) Say it once more.
\end{tcolorbox}
\section[\texorpdfstring{\gu{જકાત} (jakāt)}{જકાત (jakāt)}]{\gu{જકાત} (jakāt) — a customs duty; a toll on goods}

\label{lesson:GU-C406-jakat}



\begin{quote}
Before the new one: say the Gujarati for a visa, then the Gujarati for a boundary.
\end{quote}

\subsection*{You'll want to know: \gu{જકાત}}
\textbf{\gu{જકાત}} — \emph{jakāt} — \textquotedblleft{}a customs duty; a toll on goods\textquotedblright{}.

\begin{quote}
\gu{અમારે} \gu{સીમા} \gu{પર} \gu{જકાત} \gu{ભરવી} \gu{પડી}.
\end{quote}

\begin{quote}
\emph{amāre sīmā par jakāt bharvī paḍī.}
\end{quote}

\begin{quote}
We had to pay customs duty at the border.
\end{quote}

A feminine noun.

\begin{cousinweb}
\textbf{\gu{જકાત}} (\emph{jakāt}) is usually traced to Arabic \emph{zakāt}, \textquotedblleft{}alms, a religious levy\textquotedblright{}, which came through Persian; in Gujarati it became a duty on goods carried across a border or into a town.
\end{cousinweb}

\begin{grammarlens}[title={What you've built}]
One more word for getting around.
\end{grammarlens}

\subsection*{Your turn}
\begin{itemize}
\raggedright
  \item \emph{Say it:} \emph{jakāt}
  \item \emph{Say it:} \emph{jakāt}, once more
  \item \emph{Say it:} \emph{amāre sīmā par jakāt bharvī paḍī.} — We had to pay customs duty at the border.
  \item \emph{Recall:} say the Gujarati for a visa, then the Gujarati for a boundary, then say \emph{jakāt} again
\end{itemize}

\begin{tcolorbox}[breakable,colback=teal!4,colframe=teal!35!black,title={Before you move on}]
What does \gu{જકાત} mean? (\textquotedblleft{}A customs duty; a toll on goods\textquotedblright{}.) Say it once more.
\end{tcolorbox}
\section[\texorpdfstring{\gu{ફી} (phī)}{ફી (phī)}]{\gu{ફી} (phī) — an official fee; a charge}

\label{lesson:GU-C406-phi}



\begin{quote}
Before the new one: say the Gujarati for a boundary, then the Gujarati for a customs duty.
\end{quote}

\subsection*{You'll want to know: \gu{ફી}}
\textbf{\gu{ફી}} — \emph{phī} — \textquotedblleft{}an official fee; a charge\textquotedblright{}.

\begin{quote}
\gu{વિઝાની} \gu{ફી} \gu{ઘણી} \gu{વધારે} \gu{છે}.
\end{quote}

\begin{quote}
\emph{vijhānī phī ghaṇī vadhāre che.}
\end{quote}

\begin{quote}
The visa fee is very high.
\end{quote}

A feminine noun.

\begin{cousinweb}
\textbf{\gu{ફી}} (\emph{phī}) is borrowed from English \emph{fee}.
\end{cousinweb}

\begin{grammarlens}[title={What you've built}]
One more word for getting around.
\end{grammarlens}

\subsection*{Your turn}
\begin{itemize}
\raggedright
  \item \emph{Say it:} \emph{phī}
  \item \emph{Say it:} \emph{phī}, once more
  \item \emph{Say it:} \emph{vijhānī phī ghaṇī vadhāre che.} — The visa fee is very high.
  \item \emph{Recall:} say the Gujarati for a boundary, then the Gujarati for a customs duty, then say \emph{phī} again
\end{itemize}

\begin{tcolorbox}[breakable,colback=teal!4,colframe=teal!35!black,title={Before you move on}]
What does \gu{ફી} mean? (\textquotedblleft{}An official fee; a charge\textquotedblright{}.) Say it once more.
\end{tcolorbox}
\section[\texorpdfstring{\gu{ઠાલવવું} (ṭhālavvũ)}{ઠાલવવું (ṭhālavvũ)}]{\gu{ઠાલવવું} (ṭhālavvũ) — to empty out; to tip out}

\label{lesson:GU-C406-thalavvu}



\begin{quote}
Before the new one: say the Gujarati for a customs duty, then the Gujarati for an official fee.
\end{quote}

\subsection*{You'll want to know: \gu{ઠાલવવું}}
\textbf{\gu{ઠાલવવું}} — \emph{ṭhālavvũ} — \textquotedblleft{}to empty out; to tip out\textquotedblright{}.

\begin{quote}
\gu{તેણે} \gu{થેલી} \gu{ટેબલ} \gu{પર} \gu{ઠાલવી}.
\end{quote}

\begin{quote}
\emph{teṇe thelī ṭebal par ṭhālvī.}
\end{quote}

\begin{quote}
He emptied the bag out onto the table.
\end{quote}

Past, with \gu{તેણે} (he or she …): \gu{તેણે} \gu{થેલી} \gu{ઠાલવી}, \emph{teṇe thelī ṭhālvī}.

\begin{cousinweb}
\textbf{\gu{ઠાલવવું}} (\emph{ṭhālavvũ}) is built on \gu{ઠાલું}, \textquotedblleft{}empty\textquotedblright{}, which CDIAL lists under \emph{*ṭhalla}, \textquotedblleft{}empty\textquotedblright{}, a form it reconstructs (CDIAL 5497): to make a bag or a pot empty.
\end{cousinweb}

\begin{grammarlens}[title={What you've built}]
One more word for getting around.
\end{grammarlens}

\subsection*{Your turn}
\begin{itemize}
\raggedright
  \item \emph{Say it:} \emph{ṭhālavvũ}
  \item \emph{Say it:} \emph{ṭhālavvũ}, once more
  \item \emph{Say it:} \emph{teṇe thelī ṭebal par ṭhālvī.} — He emptied the bag out onto the table.
  \item \emph{Recall:} say the Gujarati for a customs duty, then the Gujarati for an official fee, then say \emph{ṭhālavvũ} again
\end{itemize}

\begin{tcolorbox}[breakable,colback=teal!4,colframe=teal!35!black,title={Before you move on}]
What does \gu{ઠાલવવું} mean? (\textquotedblleft{}To empty out; to tip out\textquotedblright{}.) Say it once more.
\end{tcolorbox}
